\section{Discussion and Limitations}
\subsection{What the comparison shows}
The strongest result comes from the CNN. Its local filters and pooling provide
a natural way to process the spatial structure of clothing images.
The class-level results show gains across all categories, including the
similar upper-body garments that trouble the baselines. This is consistent
with the usefulness of local spatial features in this task.

That explanation remains an interpretation. The CNN also has the largest
parameter count, and no ablation isolates the contribution of convolution,
pooling, width or regularization. The experiment demonstrates a better result
for this configuration, rather than proving a single cause for the improvement.

The MLP is a useful practical baseline. It improves substantially over Linear
and nearly matches the LSTM while requiring less recorded forward time.
The LSTM achieves similar overall accuracy with fewer parameters, but reading
the image sequentially does not provide an accuracy advantage in this run.
Its final hidden state must summarize the rows before classification.

The Transformer has the fewest parameters among the four nonlinear models,
but it has the longest recorded forward time and a lower test score than MLP,
CNN or LSTM. Attention is not automatically an improvement under a small,
fixed training configuration. Since its validation loss is still decreasing
at epoch 10, longer training is a plausible follow-up, not an established remedy.

\subsection{Limits of the evidence}
\begin{enumerate}
\item \textbf{One seed per model.} The report does not estimate variation from
initialization or data ordering. Small score differences, especially between
MLP and LSTM, should not be treated as stable.
\item \textbf{Unequal capacity and computation.} The models share an epoch budget
but have different parameter counts, operations and training durations.
The study is not a capacity-matched or compute-matched ablation.
\item \textbf{Restricted optimization.} There is no systematic learning-rate,
width, dropout or schedule search. A shared setting may suit some models better
than others, and ten epochs need not be sufficient for every configuration.
\item \textbf{Hardware and timing scope.} All cost measurements use CUDA on
the GTX 1660 SUPER. Baseline costs come from separate repetitions; classification
scores retain the original checkpoints. Inference uses only one batch size and a short benchmark,
without repeated-session timing variation. Peak model memory is not recorded.
\item \textbf{Limited image domain.} Fashion-MNIST contains small grayscale
images. The experiment does not measure robustness to color, backgrounds,
occlusion, new product styles or distribution shift.
\item \textbf{Post-evaluation analysis.} Confusions and example images explain
the existing results. They do not justify tuning against the test set and
reporting the same set as an untouched final evaluation.
\end{enumerate}

\subsection{Next experiments}
The first priority is to repeat each configuration across several seeds and
report the mean and variation of accuracy and macro-F1. Training-speed
comparisons should repeat all five models in controlled sessions on the same
hardware and timing scope to quantify runtime variation.

Next, use validation data to compare longer schedules and modest learning-rate
changes, especially for Transformer. Separate ablations can examine CNN width,
LSTM readout and patch size while controlling one factor at a time.
Augmentation and calibration are further candidates. Any new settings should
be selected through validation before a clearly documented final evaluation.
